% Nouveau EGA V (cote 25): archive cote 25, site page 14 (PDF page 15).
% typescript. Draft compared: none.
% AI-transcribed by Claude Opus 5.5 (claude-opus-5-5, adaptive thinking, effort high), Message Batches API, batch ega25b (msgbatch_01A7n5YuXubUPNK1zsepi8RG), 2026-11-13; original AI transcription, subsequently checked against both supplied scan parts; see PAGES_CHECKED.tsv and CHECK_DECISIONS.tsv.
\scanpage{25/14}
\mspage{\hand{11}}

ceci, soit Z l'ensemble des points en lesquels X n'est pas lisse sur
k, son image inverse \var{\inserted[typed]{$Z_{\hand{\eta}}$ dans $Y_{\hand{\eta}}$}}{$Z_{\hand{\eta}}$ dans $Y_{\hand{\eta}}$}{insertion}{ega25_25_014-v0001} est donc (en vertu de 2.12.)   l'ensemble des
points de $Y_{\hand{\eta}}$ en lesquels $Y_{\hand{\eta}}$ n'estb pas lisse sur k($\hand{\eta}$). Or  la codimen-%
sion de  Z dans X en x  est égale à celle de  $Z_{\hand{\eta}}$ dans $Y_{\hand{\eta}}$  en y par
raison de platitude (réf par.6). Donc l'une est \hand{$>$}k  sss l'autre l'est,
ce qui achève la démonstration.
\uline{Corollaire} 2.15.   Soit Z l'ensemble des points de X en lesquels
X n'est pas lisse sur k (resp. n'est pas  géom $R_k$ sur k, resp. n'est
\var{\margin{\unsure{//}}}{\placedmargin{\unsure{//}}}{marginal note placed}{ega25_25_014-v0002} pas séparable sur \unsure{R|k} , resp. n'est pas \var{\inserted[typed]{géom}}{géom}{insertion}{ega25_25_014-v0003} normal sur k). Pour que $Y_{\hand{\eta}}$ soit
lisse (resp.  géom $R_k$, resp séparable, resp géom normal) sur k($\hand{\eta}$ ), il
f et s  que Z soit fini.

\var{\margin{$|$}}{\placedmargin{$|$}}{marginal note placed}{ega25_25_014-v0004} Du point da vue rédaction, 2.14. et 2.15.  contiennent 2.1\unsure{2|3}. et
2.13. (qu'on pourrait donc se dispenser d'énoncer séparément), d'autre
part le corollaire est plus important pratiquement que le théorème,
qu'on pourrait donner en proposition ou lemme préliminaire, pour glori-%
fier plutôt le corollaire.

Nous pouvons  moduler \var{\xout{\unsure{xxxxparxi}}}{}{typewriter strike-out}{ega25_25_014-v0005}  le cas de la propriété (iii):
\uline{Corollaire} 2.16.  (Supposant toutjours f une immersion \var{\inserted[typed]{et de plus F cohérent}}{et de plus F cohérent}{insertion}{ega25_25_014-v0006}), sous les condi-%
\var{\margin{\unsure{//}}}{\placedmargin{\unsure{//}}}{marginal note placed}{ega25_25_014-v0007} tions de 2.7. , pour que $Z_i$ soit non immergé, il f et s que $Z_i$  le soit.
Lorsqu'il en est ainsi, la mu\hand{l}tiplicité radicielle de F en $Z_i$ \var{\xout{(ref)}}{}{typewriter strike-out}{ega25_25_014-v0008}\var{\deleted{rel}}{}{deletion}{ega25_25_014-v0009} \var{\inserted{rel.}}{rel.}{insertion}{ega25_25_014-v0010}
à k (réf) est égale à celle de $G_{\hand{\eta}}$  en $Z_i$  rel à k($\hand{\eta}$).

La première assertion est contenu dans la dernière assertion de
2.7. Pour la seconde, comme  $Y_{\hand{\eta}} \hand{\longrightarrow} X$ est plat donc le foncteur F \hand{$\rightsquigarrow$} $G_{\hand{\eta}}$
exact,  on est ramené, par restriction à un voisinage du pt gén de $Z_i$
et dévissage, au cas où F = $\underline{O}_{Z_i}$ \var{\inserted[typed]{et on peut supposer $Z_i = X$.}}{et on peut supposer $Z_i = X$.}{insertion}{ega25_25_014-v0011} . D'ailleurs on peut commencer par
se réduire au cas k alg clos, \var{\inserted[typed]{alors X est séparable sur k.}}{alors X est séparable sur k.}{insertion}{ega25_25_014-v0012} Alors \var{\xout{\unsure{$\underline{O}$xxxxx$\underline{O}_{Z_i}$xxxxx}}}{}{typewriter strike-out}{ega25_25_014-v0013} l'assertion  est
contenue dans 2.15. (iii).   On en conclut comme d'hab itude:

\var{\margin{14}}{\placedmargin{14}}{marginal note placed}{ega25_25_014-v0014}
